
This work presents an FMD-conditioned framework for evaluating formal repair methods applied to LLM-generated code. The framework measures which failure modes are present before applying methods conditioned on those failure modes.

Applied to CodeLlama-7B on HumanEval, the measurement reveals that 97.5\% of failures are syntax-level, 0\% are type-level, and 33\% of problems are Z3-eligible. SynCode directionally reduces AST failures (delta_ast = 0.075, consistent across two experiments, not statistically significant at N = 20). mypy is operational but produces no type error feedback for annotation-free LLM code. Z3 eligibility scope is established but repair effectiveness is untested.

The three-method pipeline reduces to a two-method pipeline for this model-benchmark pair. The SynCode-Z3 complementarity test (h-m3), the most theoretically motivated remaining experiment, was not executed in this pipeline run. Extensions include: (1) full N = 164 SynCode run for statistical confirmation; (2) execution of h-m3 for SynCode-Z3 C_score measurement; (3) evaluation with annotation-prompted or larger models to test whether the zero type stratum persists; (4) Z3 eligibility heuristic unification across experiments.

The principal empirical contribution is the demonstration that formal repair channel scope is measurable and model-specific. For CodeLlama-7B on HumanEval, the type-checking repair channel is structurally inapplicable -- a finding that is not predictable from theoretical analysis and that was not previously measured for this model-benchmark combination.
